% Part: first-order-logic
% Chapter: introduction
% Section: formulas

\documentclass[../../../include/open-logic-section]{subfiles}

\begin{document}

\olfileid[tr]{fol}{int}{fml}

\section{\usetoken{P}{formula}}

Birinci dereceden mantığın !!{sentence}s ifadelerini titizlikle
belirlemek ve !!{variable}s kullanımından doğan sorunları ele almak
için şu yaklaşımı kullanacağız. Önce \emph{farklı} bir ifadeler kümesi
tanımlarız: !!{formula}s. Bunu yaptıktan sonra !!{variable}s
ifadelerinin bunlarda oynadığı rolü ele alabilir ve ``serbest'' ile
``bağlı'' !!{variable}s fikirlerine dayanarak !!a{sentence} ifadesini
(yani serbest !!{variable}s içermeyen !!a{formula} ifadesini)
tanımlayabiliriz. Bunu yalnızca ``!!{sentence}'' tanımını daha kolay
yönetilebilir kıldığı için değil, semantik sağlama kavramını tanımlama
biçimimiz açısından belirleyici olacağı için de yaparız.

``!!{formula}''yı yalnızca tek bir !!{predicate}~$\Obj P$, tek bir
!!{constant}~$\Obj a$ ve yalnızca $\lnot$, $\land$ ve~$\lexists$
mantıksal simgelerini içeren basit bir birinci dereceden dil için
tanımlayalım. Tam tanımlarımız çok daha genel olacaktır: sonsuz sayıda
!!{predicate}s ve !!{constant}s bulunmasına izin vereceğiz. Hatta
!!{constant}s ve !!{variable}s ile birleştirilerek ``terimler''
oluşturabilen !!{function}s ifadelerini de ele alacağız. Şimdilik
yalnızca $\Obj a$ ile değişkenler terimimiz olacak. Sonsuz sayıda
!!{variable}s gerekir. Resmî olarak $\Obj v_0$, $\Obj v_1$, \dots\
simgelerini değişken olarak kullanacağız.

\begin{defn}
\emph{!!{formula}s} kümesi~$\Frm$ aşağıdaki gibi tanımlanır:
\begin{enumerate}
\item\ollabel{fmls-atom} $\Atom{\Obj P}{\Obj a}$ ve
  $\Atom{\Obj P}{\Obj v_i}$ birer !!{formula}s ifadesidir
  ($i\in\Nat$).

\tagitem{prvNot}{\ollabel{fmls-not}$!A$ !!a{formula} ise
  $\lnot !A$ da bir !!{formula} ifadesidir.}{}

\tagitem{prvAnd}{$!A$ ve $!B$ !!{formula}s ise $(!A\land !B)$
  bir !!a{formula} ifadesidir.}{}

\tagitem{prvEx}{\ollabel{fmls-ex}$!A$ !!a{formula} ve $x$
  !!a{variable} ise $\lexists[x][!A]$ bir !!a{formula} ifadesidir.}{}

\tagitem{limitClause}{\ollabel{fmls-limit}Bunların dışında hiçbir şey
  !!a{formula} değildir.}{}
\end{enumerate}
\end{defn}

\olref{fmls-atom}, herhangi bir $i\in\Nat$ için
$\Atom{\Obj P}{\Obj a}$ ile $\Atom{\Obj P}{\Obj v_i}$ ifadelerinin
!!{formula}s olduğunu söyler. Bunlara \emph{atomik} !!{formula}s
denir. Başlangıç noktamızı bunlar verir. Öteki maddeler, daha önce
oluşturduğumuz !!{formula}s ifadelerinden yenilerini oluşturma yolları
verir. Örneğin $\Atom{\Obj P}{\Obj v_2}$,~\olref{fmls-atom} uyarınca
zaten !!a{formula} olduğundan \olref{fmls-not} uyarınca
$\lnot\Atom{\Obj P}{\Obj v_2}$ bir !!a{formula} ifadesidir. Ardından
\olref{fmls-ex} uyarınca
$\lexists[\Obj v_2][\lnot\Atom{\Obj P}{\Obj v_2}]$ başka bir
!!{formula} ifadesidir; süreç böyle sürer. \olref{fmls-limit}, \emph{yalnızca}
bu biçimde oluşturabileceğimiz dizgelerin !!{formula}s sayıldığını
söyler. Özellikle $\lexists[\Obj v_0][\Atom{\Obj P}{\Obj a}]$ ile
$\lexists[\Obj v_0][\lexists[\Obj v_0][\Atom{\Obj P}{\Obj a}]]$
!!{formula}s \emph{sayılır}; fakat gereksiz dış parantezler nedeniyle
$(\lnot\Atom{\Obj P}{\Obj a})$ sayılmaz.

!!{formula}s ifadelerini bu biçimde tanımlamaya \emph{tümevarımlı
tanım} denir ve bu tanım, \emph{yapısal tümevarım} denen bir tümevarımla
ispat biçimini kullanarak !!{formula}s hakkında ispatlar vermemizi
sağlar. Bunlar \olref[mth][ind][idf]{sec} ve
\olref[mth][ind][sti]{sec} içinde genel olarak ele alınmıştır; daha
sonraki ispatlara geçmeden önce bu kısımları gözden geçirmelisiniz.
Temel fikir şudur: Bir şeyin bütün !!{formula}s ifadeleri için doğru
olduğunu ispatlamak istiyorsanız önce atomik !!{formula}s için doğru
olduğunu, ardından herhangi bir~$!A$ (ve~$!B$) !!{formula} ifadesi
için doğru \emph{ise} $\lnot !A$, $(!A\land !B)$ ve
$\lexists[x][!A]$ için de doğru olduğunu gösterirsiniz. Örneğin bu,
özelliğin $\lexists[\Obj v_2][\lnot\Atom{\Obj P}{\Obj v_2}]$ için
doğru olduğunu ispatlar: ilk bölümden atomik
!!{formula}~$\Atom{\Obj P}{\Obj v_2}$ için doğru olduğunu bilirsiniz.
İkinci bölümle $\lnot\Atom{\Obj P}{\Obj v_2}$ için ve yine aynı
bölümle $\lexists[\Obj v_2][\lnot\Atom{\Obj P}{\Obj v_2}]$ için doğru
olduğunu elde edersiniz. Bütün !!{formula}s, atomik !!{formula}s
ifadelerinden tümevarımlı olarak üretildiğinden yöntem bunların hepsi
için çalışır.

\end{document}
